% ARI5002 Course Project Report — IEEE Conference Format
% Compile with: pdflatex main.tex (twice), then bibtex main, then pdflatex twice more

\documentclass[conference]{IEEEtran}
\IEEEoverridecommandlockouts

\usepackage{cite}
\usepackage{amsmath,amssymb,amsfonts}
\usepackage{algorithm}
\usepackage{algorithmic}
\usepackage{graphicx}
\usepackage{textcomp}
\usepackage{xcolor}
\usepackage{booktabs}
\usepackage{multirow}
\usepackage{url}
\usepackage{bm}

\def\BibTeX{{\rm B\kern-.05em{\sc i\kern-.025em b}\kern-.08em
    T\kern-.1667em\lower.7ex\hbox{E}\kern-.125emX}}

\begin{document}

\title{Replication and Extension of POMO for Heterogeneous\\
Electric-Diesel Fleet Routing with Time Windows}

\author{
  \IEEEauthorblockN{Egemen Eroglu}
  \IEEEauthorblockA{
    \textit{YOUR UNIVERSITY NAME}\\
    erogluegemen@gmail.com
  }
}

\maketitle

%% ─────────────────────────────────────────────────────────────────────────────
\begin{abstract}
This paper presents a replication and extension of POMO (Policy Optimization
with Multiple Optima)~\cite{kwon2020pomo}, a deep reinforcement learning
framework for combinatorial optimization. We faithfully re-implement the
POMO training algorithm and Transformer-based encoder-decoder architecture
in PyTorch, then extend it to three additional problem variants: the Vehicle
Routing Problem with Time Windows (VRPTW), an electric vehicle (EV) range
constraint, and a heterogeneous fleet via vehicle-type embeddings. On
synthetic CVRP benchmarks our re-implementation achieves mean tour lengths
of 4.327 (n=20) and 8.922 (n=50) with 8$\times$ coordinate augmentation,
outperforming OR-Tools (15\,s limit) by 10.75\% and 0\% respectively.
On synthetic VRPTW the model beats OR-Tools by 9.14\% on n=20 instances.
We further apply the trained model to real operational data from DHL Istanbul
(905 routes across three depots with a mixed diesel/EV fleet), analysing the
substantial distribution shift that limits direct transfer. On select
small-scale instances the model reduces driven distance by 41--54\% relative
to actual driver routes.
\end{abstract}

\begin{IEEEkeywords}
deep reinforcement learning, combinatorial optimization, vehicle routing,
POMO, VRPTW, electric vehicles, replication study
\end{IEEEkeywords}

%% ─────────────────────────────────────────────────────────────────────────────
\section{Introduction}

Combinatorial optimization underlies a large class of real-world logistics
problems. The Capacitated Vehicle Routing Problem (CVRP) and its extension
with time windows (VRPTW) are canonical NP-hard problems that underpin
last-mile delivery planning~\cite{dantzig1959truck}. Classical exact solvers
such as LKH3~\cite{helsgott2017lkh3} and OR-Tools~\cite{ortools} can find
near-optimal solutions but require minutes to hours per instance, making
them unsuitable for online re-routing or large-scale fleet management.

Deep reinforcement learning (DRL) has emerged as a promising alternative.
Construction-type neural methods~\cite{vinyals2015pointer, kool2019attention}
learn to build solutions autoregressively, producing near-optimal tours in
milliseconds at inference time. POMO~\cite{kwon2020pomo}, introduced at
NeurIPS~2020, advances this line by exploiting the fact that every customer
node is an equally valid starting point for a VRP solution. By training with
$n$ simultaneous rollouts---one per starting node---and using their mean
return as a low-variance baseline, POMO achieves faster convergence and
stronger solutions than prior construction methods.

\textbf{Contributions.} This work (i) re-implements POMO from first
principles in PyTorch and reproduces its core CVRP results; (ii) extends the
model to VRPTW with time-feasibility masking and to EVs with a cumulative
range constraint; (iii) introduces vehicle-type embeddings to support
heterogeneous diesel/EV fleets; (iv) applies the trained model to real DHL
Istanbul operational data and quantifies the distribution shift between the
synthetic training distribution and real last-mile delivery instances.

%% ─────────────────────────────────────────────────────────────────────────────
\section{Background and Related Work}

\subsection{Neural Construction Methods}

Pointer Networks~\cite{vinyals2015pointer} were among the first deep-learning
models to solve routing problems by outputting permutations autoregressively
via a pointer mechanism. The Attention Model (AM)~\cite{kool2019attention}
replaces recurrence with a Transformer encoder and a single-layer attention
decoder, trained with REINFORCE~\cite{williams1992} using a greedy rollout
baseline. POMO~\cite{kwon2020pomo} builds on AM by initializing $N$ parallel
rollouts from different starting nodes per instance, using the batch mean of
these returns as a shared baseline. This approach both reduces gradient
variance and prevents the policy from anchoring to a single starting node.

\subsection{VRP Variants}

The VRPTW~\cite{solomon1987algorithms} augments CVRP with per-customer hard
time windows $[e_i, l_i]$: a vehicle arriving before $e_i$ must wait, and
arriving after $l_i$ is infeasible. The Heterogeneous Fleet VRP (HFVRP)
introduces vehicles of different capacities and costs; here we additionally
distinguish electric vehicles (EVs) that carry a maximum driving range
constraint.

%% ─────────────────────────────────────────────────────────────────────────────
\section{Problem Formulation}

\subsection{Capacitated VRP}

Given a depot node $v_0$ and $n$ customer nodes $\{v_1,\ldots,v_n\}$ with
coordinates $\mathbf{x}_i \in [0,1]^2$ and demands $d_i > 0$, and a vehicle
of capacity $C$, the CVRP seeks a minimum-length set of routes starting and
ending at the depot such that (a) every customer is visited exactly once and
(b) the total demand on each route does not exceed $C$. Multiple depot visits
by the same vehicle (multi-trip) are permitted; the objective is the total
circuit length.

\subsection{VRPTW Extension}

In VRPTW each customer $v_i$ additionally carries a time window
$[e_i, l_i]$. Let $t_i$ denote the vehicle's arrival time at $v_i$ and
$s$ the service time. Feasibility requires $t_i \le l_i$ for all $i$.
Waiting is allowed: $\max(t_i, e_i)$ is the actual service start. The node
features are augmented from 3-d to 5-d:
$\mathbf{f}_i = (x_i,\, y_i,\, d_i/C,\, e_i/T,\, l_i/T)$
where $T$ is the time horizon.

\subsection{Heterogeneous Fleet and EV Range}

We model a heterogeneous fleet where vehicle type $k \in
\{\text{SCV}, \text{MCV}, \text{LCV}, \text{LCV\_BEV}\}$ is fixed per route.
LCV\_BEV vehicles additionally carry a range limit $R$: the accumulated
route distance before the next return to depot may not exceed $R$. This is
enforced at each decoding step by masking candidates $v_j$ for which
$d_{\text{cumul}} + \text{dist}(v_{\text{cur}}, v_j) +
\text{dist}(v_j, v_0) > R$.

%% ─────────────────────────────────────────────────────────────────────────────
\section{POMO: Architecture and Training}

\subsection{Encoder}

The encoder maps node features $\mathbf{F} \in \mathbb{R}^{(n+1)\times f}$
($f=3$ for CVRP, $f=5$ for VRPTW) to embeddings
$\mathbf{H} \in \mathbb{R}^{(n+1)\times d}$ via a linear projection followed
by $L$ Transformer layers with Pre-LayerNorm:
\begin{equation}
  \mathbf{H}^{(l)} = \text{FFN}\!\left(\text{MHA}\!\left(\mathbf{H}^{(l-1)}\right)\right),
\end{equation}
where MHA denotes multi-head self-attention~\cite{vaswani2017attention} and
FFN a two-layer feed-forward network with ReLU activation. Hyperparameters
follow the original paper: $d=128$, 8 attention heads, $d_\text{ff}=512$,
$L=6$.

A graph-level embedding is computed as the mean of node embeddings:
$\bar{\mathbf{h}} = \frac{1}{n+1}\sum_{i=0}^n \mathbf{H}_i$.

\subsection{Decoder}

At each step $t$ the decoder assembles a context vector
\begin{equation}
  \mathbf{q}_t = \text{Linear}\!\left(\left[\bar{\mathbf{h}},\,
  \mathbf{H}_{v_t},\, r_t / C\right]\right),
\end{equation}
where $v_t$ is the current node, $r_t$ is the remaining capacity, and
$[\cdot,\cdot,\cdot]$ denotes concatenation. For VRPTW the current
normalised time $\tau_t / T$ is appended. For the heterogeneous fleet a
vehicle-type embedding (4-class, 16-d, projected to $d$) is further
concatenated.

Single-head attention computes compatibility scores
\begin{equation}
  u_{ti} = C_{\text{clip}} \cdot \tanh\!\left(
    \frac{\mathbf{q}_t \mathbf{H}_i^\top}{\sqrt{d}}
  \right),
\end{equation}
with $C_{\text{clip}}=10$ following~\cite{kwon2020pomo}. After applying the
infeasibility mask $\mathbf{m}_t$ (visited nodes, capacity, time windows,
EV range), the action is sampled from
$\pi_\theta(a_t \mid s_t) = \text{softmax}(\mathbf{u}_t - \mathbf{m}_t \cdot
\infty)$.

\subsection{POMO Training}

POMO exploits the $n$ symmetries of the VRP by running one rollout per
customer node as starting point. For batch size $B$ and $N=n$ starts, the
encoder runs once per instance and is shared across all $N$ rollouts. The
gradient estimator is
\begin{equation}
  \nabla_\theta J(\theta) \approx \frac{1}{BN}
  \sum_{i=1}^{B}\sum_{j=1}^{N}
  \bigl(R(\bm{\tau}_i^j) - b_i\bigr)\,
  \nabla_\theta \log p_\theta(\bm{\tau}_i^j),
\end{equation}
where $b_i = \frac{1}{N}\sum_{j=1}^N R(\bm{\tau}_i^j)$ is the
within-instance shared baseline~\cite{kwon2020pomo}. This baseline has zero
mean for the POMO rollouts, reducing gradient variance without a separately
trained critic.

\subsection{Inference}

At inference, $n$ greedy (argmax) rollouts are executed in parallel and the
shortest tour selected. For 8$\times$ augmentation, the input coordinates
are transformed by the 8 dihedral symmetries of the unit square (4 rotations
$\times$ 2 reflections) and all $8n$ rollouts are scored; the global minimum
is returned. Algorithm~\ref{alg:inference} summarises this procedure.

\begin{algorithm}[t]
\caption{POMO Inference with Augmentation}
\label{alg:inference}
\begin{algorithmic}[1]
\REQUIRE instance $s$, policy $\pi_\theta$, augmentations $K=8$
\STATE $\tau^* \leftarrow \varnothing$,\quad $R^* \leftarrow -\infty$
\FOR{$k = 1$ \textbf{to} $K$}
  \STATE $s_k \leftarrow \textsc{Augment}(s, k)$
  \STATE $\mathbf{H}_k \leftarrow \textsc{Encoder}(s_k)$
  \FOR{$j = 1$ \textbf{to} $N$}
    \STATE $\bm{\tau}_k^j \leftarrow \textsc{GreedyRollout}(\mathbf{H}_k, j, \pi_\theta)$
  \ENDFOR
\ENDFOR
\RETURN $\arg\max_{k,j} R(\bm{\tau}_k^j)$
\end{algorithmic}
\end{algorithm}

%% ─────────────────────────────────────────────────────────────────────────────
\section{Extensions}

\subsection{VRPTW Time-Window Masking}

At each decoding step, the arrival time at candidate $v_j$ is estimated as
$t_j = t_{\text{cur}} + \text{dist}(v_{\text{cur}}, v_j) / v$, where $v$ is
vehicle speed. Node $v_j$ is masked if $t_j > l_j$ (hard deadline). After
completing a trip, the vehicle returns to the depot and time resets to the
arrival instant. Node features include normalised time windows
$(e_i/T,\, l_i/T)$, and the current normalised time $\tau_t/T$ is provided
to the decoder at each step.

\subsection{EV Range Constraint}

For electric vehicles the cumulative distance since the last depot visit is
tracked. Candidate $v_j$ is masked if
$d_{\text{cumul}} + \text{dist}(v_{\text{cur}},v_j) +
\text{dist}(v_j, v_0) > R$, ensuring the vehicle can always return to the
depot to recharge. When the vehicle returns to the depot, $d_{\text{cumul}}$
resets to zero.

\subsection{Vehicle-Type Embedding}

To support heterogeneous fleets, a learnable embedding
$\mathbf{e}_k \in \mathbb{R}^{16}$ is defined for each vehicle type
$k \in \{\text{SCV}, \text{MCV}, \text{LCV}, \text{LCV\_BEV}\}$.
At decode time this is concatenated to the context vector before the
compatibility computation, and a linear layer projects the extended context
back to dimension $d$.

%% ─────────────────────────────────────────────────────────────────────────────
\section{Experimental Setup}

\subsection{Implementation}

All code is written in PyTorch and runs device-agnostically on CUDA, Apple
MPS, or CPU. Training is performed on Google Colab (GPU). The encoder-decoder
uses the hyperparameters described in Section~IV: $d=128$, 8 heads,
$d_\text{ff}=512$, $L=6$, $C_\text{clip}=10$. We train with Adam~\cite{kingma2015adam}
at learning rate $10^{-4}$, batch size 64 (2560 synthetic instances per epoch),
for 100 epochs on synthetic instances generated on the fly.

Two models are trained separately:
\begin{itemize}
  \item \textbf{CVRP model} ($f=3$): node features $(x, y, d/C)$.
  \item \textbf{VRPTW model} ($f=5$): node features $(x, y, d/C, e/T, l/T)$.
\end{itemize}

\subsection{Synthetic Benchmark Instances}

For CVRP, instances follow the standard setup: $n \in \{20, 50\}$ customer
nodes, coordinates $\sim \mathcal{U}[0,1]^2$, demands $\sim
\mathcal{U}\{1,\ldots,9\}$, vehicle capacity $C=50$. We evaluate on 100
held-out instances per problem size.

\textbf{Remark on capacity.}
The original POMO paper~\cite{kwon2020pomo} follows Kool
et al.~\cite{kool2019attention} and uses $C=30$ for $n=20$ and $C=40$ for
$n=50$. Our implementation uses $C=50$ uniformly (the setting from
configs/default.yaml, matching the DHL fleet planning context). A lower
capacity forces more depot returns and therefore longer tours; this causes
our absolute tour lengths to be shorter than the original paper's Table~3
but the relative trends are preserved.

For VRPTW, time windows are generated with half-width $h \sim
\mathcal{U}[0.15T, 0.35T]$ centred within $[0, T]$, with $T=480$ minutes.
Since travel times are Euclidean distances in $[0,1]^2$ (max $\approx 1.41$),
all time windows are effectively unconstrained at the scale of $T=480$; this
is a known limitation discussed in Section~VIII.

\subsection{Baselines}

We compare against:
\begin{itemize}
  \item \textbf{Nearest Neighbour (NN)}: greedy construction from the depot,
    repeatedly visiting the closest unvisited feasible customer.
  \item \textbf{OR-Tools}~\cite{ortools}: Google's constraint-programming
    solver with a 15-second time limit per instance, using $k =
    \lceil D_\text{total}/C \rceil$ vehicles to model multi-trip behaviour.
\end{itemize}

All baselines and POMO are evaluated on the same 100 instances.

\subsection{Real-World Dataset}

We additionally evaluate on operational data from DHL's Istanbul
last-mile network, covering January 2026 across three depots: Sabiha
Gök\c{c}en (SAW), Istanbul Airport (IGA), and a Central European Transfer
(CET) hub. From the raw delivery records we extract 905 VRP instances,
grouped by depot, route code, and date. Each instance comprises a depot, a
set of customer stops with GPS coordinates, demand (kg), actual time
windows, and a vehicle assignment (type and capacity from the fleet
registry). The fleet includes SCV, MCV, and LCV diesel vans plus LCV\_BEV
electric vehicles (Ford e-Transit).

Haversine distances are pre-computed and stored per instance. Coordinates
are min-max normalised to $[0,1]^2$ for model input. The DHL driver route
is reconstructed by sorting stops by recorded arrival time, serving as the
baseline.

%% ─────────────────────────────────────────────────────────────────────────────
\section{Results}

\subsection{Synthetic CVRP and VRPTW Benchmark}

Table~\ref{tab:benchmark} reports mean tour lengths on 100 held-out instances
for each method and problem size. POMO with 8$\times$ augmentation
outperforms OR-Tools (15\,s) on CVRP n=20 by \textbf{10.75\%} and on VRPTW
n=20 by \textbf{9.14\%}, while producing results in under one second.
The nearest-neighbour heuristic is 27\% worse than POMO on CVRP n=20,
confirming POMO's superiority over classical greedy construction.

\begin{table}[t]
\caption{Synthetic Benchmark: Mean Tour Length (100 instances, $C=50$)}
\label{tab:benchmark}
\centering
\begin{tabular}{llccc}
\toprule
Problem & $n$ & Method & Length & Gap vs.\ OR-Tools \\
\midrule
\multirow{4}{*}{CVRP}
  & \multirow{4}{*}{20}
    & Nearest Neighbour   & 6.173  & $+27.3\%$ \\
  && OR-Tools (15\,s)    & 4.848  & --- \\
  && POMO greedy (ours)  & 4.469  & $-7.81\%$ \\
  && POMO +8$\times$aug  & 4.327  & $-10.75\%$ \\
\midrule
\multirow{2}{*}{CVRP}
  & \multirow{2}{*}{50}
    & POMO greedy (ours)  & 9.139  & --- \\
  && POMO +8$\times$aug  & 8.922  & --- \\
\midrule
\multirow{3}{*}{VRPTW}
  & \multirow{3}{*}{20}
    & OR-Tools (15\,s)   & 5.039  & --- \\
  && POMO greedy (ours)  & 4.752  & $-5.71\%$ \\
  && POMO +8$\times$aug  & 4.579  & $-9.14\%$ \\
\midrule
\multirow{2}{*}{VRPTW}
  & \multirow{2}{*}{50}
    & POMO greedy (ours)  & 9.696  & --- \\
  && POMO +8$\times$aug  & 9.390  & --- \\
\bottomrule
\end{tabular}
\end{table}

\subsection{Comparison with Original Paper}

Table~\ref{tab:comparison} compares our CVRP results with those reported in
Kwon et al.~\cite{kwon2020pomo}. Because of the different capacity settings
($C=50$ vs $C=30$ for $n=20$, $C=40$ for $n=50$ in the original), absolute
tour lengths are not directly comparable. However, the key qualitative
trends are reproduced:

\begin{enumerate}
  \item Augmentation consistently improves over greedy decoding.
  \item The relative gap between greedy and augmented decoding grows with
    problem size ($n=20$: $-3.2\%$; $n=50$: $-2.4\%$ for ours vs $-0.5\%$;
    $-0.7\%$ for the original), suggesting our shorter training (100 epochs)
    leaves more room for augmentation to compensate.
  \item POMO outperforms OR-Tools, consistent with the original finding that
    POMO beats OR-Tools on CVRP20 and CVRP50.
\end{enumerate}

\begin{table}[t]
\caption{Comparison with Original POMO Results on CVRP\\
  {\small (Original: $C=30$ for $n=20$, $C=40$ for $n=50$; Ours: $C=50$)}}
\label{tab:comparison}
\centering
\begin{tabular}{lccc}
\toprule
Method & CVRP20 & CVRP50 & Augm.\ gain ($n=20$) \\
\midrule
POMO single trajec.\ (original) & 6.35 & 10.74 & ---\\
POMO no aug.\ (original)        & 6.17 & 10.49 & $-2.8\%$ \\
POMO $\times$8 aug.\ (original) & 6.14 & 10.42 & $-0.5\%$ \\
\midrule
POMO greedy (ours)               & 4.469 & 9.139 & ---\\
POMO $\times$8 aug.\ (ours)      & 4.327 & 8.922 & $-3.2\%$ \\
\bottomrule
\end{tabular}
\end{table}

\begin{figure}[t]
  \centering
  \includegraphics[width=\columnwidth]{../viz/figures/benchmark_comparison.png}
  \caption{Grouped bar chart of mean tour lengths for each method and problem
    configuration. NN: Nearest Neighbour. POMO greedy uses a single argmax
    per start; POMO 8$\times$aug applies 8 coordinate symmetry transforms.}
  \label{fig:benchmark}
\end{figure}

\subsection{EV Range Sensitivity}

Fig.~\ref{fig:ev} shows the effect of the EV range constraint on CVRP n=20
instances. Routes remain feasible and achieve the same tour length for
$R \ge 3.0$ distance units (max possible inter-node distance in $[0,1]^2$
is $\sqrt{2} \approx 1.41$, so $R=3.0$ is effectively unlimited). Only
at $R=2.0$ does the constraint become binding, increasing mean tour length
by \textbf{8.8\%} as the vehicle must return to depot more frequently to
avoid range depletion.

\begin{figure}[t]
  \centering
  \includegraphics[width=0.85\columnwidth]{../viz/figures/ev_range_sensitivity.png}
  \caption{Mean tour length vs.\ EV range constraint (CVRP, $n=20$, 100
    instances). The constraint only binds at $R=2.0$, forcing additional
    depot returns and increasing tour length by 8.8\%.}
  \label{fig:ev}
\end{figure}

\subsection{Real-World Transfer: DHL Istanbul}

Applying the CVRP model directly to the 905 real DHL instances yields
substantially worse routes than the drivers' actual sequences on average.
This confirms a significant distribution shift between the synthetic training
domain and real last-mile delivery data.

Fig.~\ref{fig:distshift} summarises four key mismatches identified by
comparing training data statistics with the real dataset:

\begin{enumerate}
  \item \textbf{Demand scale (14.7$\times$ gap)}: Training uses
    normalised demands $d_i/C \in [0.02, 0.18]$. Real packages weigh
    0.1--30\,kg against vehicle capacities of 500--1500\,kg, yielding
    $d_i/C \approx 0.0001$--$0.075$. The model has never encountered
    such small fractional demands.
  \item \textbf{Coordinate distribution}: Training is uniform random in
    $[0,1]^2$. Real coordinates are clustered in distinct urban districts
    of Istanbul after min-max normalisation.
  \item \textbf{Time window structure}: 47.9\% of real stops carry
    all-day windows ($[0, 1439]$), while the rest show narrow operational
    windows. Training windows span $[0.3T, 0.7T]$ half-width---a
    distribution not present in real data.
  \item \textbf{VRPTW unit mismatch}: The training time horizon $T=480$
    minutes is much larger than the maximum Euclidean travel time in
    $[0,1]^2$ ($\approx 1.41$ units with $v=1$). Time-window constraints
    are therefore never binding during synthetic VRPTW training, so the
    VRPTW model is functionally equivalent to a CVRP model with a 5-d
    input.
\end{enumerate}

Despite the aggregate failure, on a selected subset of 28--30 node
instances from CET depot---which are geographically compact and have few
all-day time windows---the CVRP model achieves 41--54\% reductions in
total route distance compared with actual driver sequences
(Table~\ref{tab:realworld}). Route maps for these instances are provided
as interactive HTML files in the accompanying code submission.

\begin{figure}[t]
  \centering
  \includegraphics[width=\columnwidth]{../viz/figures/distribution_shift.png}
  \caption{Distribution shift analysis. Each panel compares the training data
    statistics (blue) with the DHL Istanbul dataset (orange). Significant
    mismatches exist in demand normalisation, coordinate distribution, and
    time-window structure.}
  \label{fig:distshift}
\end{figure}

\begin{table}[t]
\caption{POMO vs.\ DHL Driver Routes on Best-Case Real Instances\\
  {\small (CVRP model, 8$\times$ augmentation)}}
\label{tab:realworld}
\centering
\begin{tabular}{llcccr}
\toprule
Depot & Route & $n$ & Type & DHL (km) & POMO (km) \\
\midrule
CET & CECD & 30 & SCV (diesel) & 19.7 & 9.1 \\
CET & CEGA & 22 & SCV (diesel) & 20.4 & 11.9 \\
IGA & ISCB & 28 & LCV\_BEV (EV) & 54.3 & 27.1 \\
SAW & EADB & 28 & SCV (diesel) & 72.7 & 38.0 \\
\bottomrule
\multicolumn{4}{l}{\small Route distance reductions: 38--54\%.} \\
\end{tabular}
\end{table}

%% ─────────────────────────────────────────────────────────────────────────────
\section{Discussion}

\subsection{Fidelity of the Replication}

Our implementation successfully reproduces the core POMO mechanisms: shared
within-instance REINFORCE baseline, multi-start greedy inference, and
$8\times$ coordinate augmentation. The benefit of augmentation over greedy
decoding ($-3.2\%$ for n=20) and the consistent superiority over OR-Tools
are both in line with the original paper's findings. The main deviation is
quantitative: our tour lengths are systematically lower because we train and
evaluate at $C=50$ uniformly, while the original paper uses $C=30$ ($n=20$)
and $C=40$ ($n=50$), which induce more depot returns and longer tours.

\subsection{VRPTW Training Limitation}

The most significant implementation gap concerns the VRPTW model. The
training configuration sets the time horizon to $T=480$ minutes (an 8-hour
delivery day), but the synthetic Euclidean distances in $[0,1]^2$ are at
most $\approx 1.41$ units. With unit speed, travel times are orders of
magnitude smaller than any time window half-width ($\ge 72$ minutes), so the
time-window feasibility mask is never triggered during training. Consequently,
the VRPTW model learns to produce valid CVRP solutions but does not learn to
respond to tight time windows. A correct fix requires scaling the travel
distances to the same units as $T$ (e.g., setting $T=2.0$ to match the
unit-square scale) and retraining.

\subsection{Real-World Distribution Shift}

The transfer failure on the 905 DHL instances is primarily driven by three
compounding mismatches. First, the 14.7$\times$ demand normalisation gap
means the capacity constraint is almost never binding in real instances,
collapsing the problem to an uncapacitated TSP for which the model was not
trained. Second, urban GPS coordinates cluster spatially in ways that are
absent from the uniform training distribution. Third, the proportion of
all-day windows (48\%) introduces a mode not seen in training.

These observations are a valuable finding in their own right: they motivate
domain adaptation techniques such as fine-tuning on real instances or mixed
training distributions, which are identified as the primary avenue for
future work.

\subsection{Computational Cost}

Training the CVRP model for 100 epochs on Google Colab required
approximately 2--3 hours on a single GPU. Inference on 100 instances with
8$\times$ augmentation completed in under 60 seconds on Apple MPS hardware,
demonstrating the practical speed advantage of POMO over OR-Tools (which
requires $\sim$1500\,s for the same 100 instances).

%% ─────────────────────────────────────────────────────────────────────────────
\section{Conclusion}

We have replicated POMO's CVRP solver and extended it to VRPTW, EV range
constraints, and heterogeneous vehicle-type embeddings. On synthetic
benchmarks our implementation surpasses OR-Tools by up to 10.75\% (CVRP)
and 9.14\% (VRPTW) with near-zero inference time. The EV range extension
correctly enforces range feasibility, incurring an 8.8\% distance penalty
only when the constraint is actively binding.

Application to real DHL Istanbul data reveals a significant distribution
shift that prevents direct transfer of the synthetically-trained model.
We identify four specific mismatch sources---demand scale, coordinate
clustering, time-window structure, and VRPTW unit mismatch in training---and
provide a quantitative analysis. The replication demonstrates that POMO's
algorithmic core (multi-start REINFORCE, shared baseline, coordinate
augmentation) is sound and reproducible, while also highlighting that
real-world deployment requires targeted domain adaptation.

Future work should address: (i) retraining VRPTW with a corrected time
scale ($T=2.0$); (ii) fine-tuning on the DHL dataset with mixed training
distributions; (iii) adding a comparison with the Attention Model
baseline~\cite{kool2019attention} on our benchmark.

%% ─────────────────────────────────────────────────────────────────────────────
\bibliographystyle{IEEEtran}
\bibliography{references}

\end{document}
